\begin{center}
  {\zihao{2}\bfseries \PaperTitle\par}
  \vspace{1.2em}
  {\zihao{3}\bfseries 摘要\par}
\end{center}
\vspace{0.8em}

本文构造一个有限容量下的合成资源分配任务，用于验证“正式运行结果到中文数学建模论文”的发布闭环，
并不提供任何竞赛问题结论。每个物品具有成本与价值，任务是在容量约束下最大化总价值。我们比较按
单位成本价值排序的密度贪心法与遍历全部可行子集的穷举法，并将实验定义、原始观测、指标、图表数据、
绘图脚本和独立检查固定在同一运行包中。

正式运行计划包含 $3$ 个合成实例，其中 $2$ 个完成、$1$ 个按设计失败，故总体结果为部分完成。
在两个完成实例上，密度贪心法的完整观测均值为 \BaselineMeanValue{} \BaselineMeanValueUnit，
穷举法均值为 \AlternativeMeanValue{} \AlternativeMeanValueUnit；这些数值由运行清单确定性生成，
不是人工抄录。逐实例结果显示，一个实例存在 $6$ 个 value 的最优性差距，另一个实例两法同值。
由于没有重复随机试验或已定义的抽样分布，本文不绘制误差区间，也不把两个完整观测推广为总体性能结论。

发布阶段进一步检查引用、匿名信息、关键数字、PDF 元数据、字体嵌入和可搜索文本层，并栅格化全部页面
供实际视觉复核。该演示说明：样式、颜色、线型和图例可在不改变 Metric 的情况下重绘；一旦改变样本、
聚合、单位或正式运行引用，则必须返回科学结果修订。

\vfill
\noindent\textbf{关键词：}\PaperKeywords
